% Research Objectives, Content, and Key Problems
\section{Research Objectives}

To achieve the stated aim and test the hypotheses, this research
is structured around four distinct, measurable objectives.

\subsection{Objective 1: To Design a Sparse-Sampling Framework for Reasoning Segmentation}
The first objective is to architect the software pipeline that allows for the decoupling of "Reasoning" and "Tracking."

\begin{itemize}
    \item \textbf{Rationale:} Current models are monolithic; they cannot easily skip frames without losing context.
          I must design a modular system where the "Semantic Module" (\cite[VideoLISA]{bai_one_2024}) and the
          "Interpolation Module" are distinct but synchronized.
    \item \textbf{Implementation: } This involves developing a Keyframe Selector Algorithm.
          Rather than selecting keyframes arbitrarily (e.g., every 5th frame), the algorithm will analyze pixel
          difference metrics (Optical Flow magnitude) to trigger the heavy Deep Learning model only when the scene
          changes significantly. This ensures that compute resources are invested intelligently.
\end{itemize}

\subsection{Objective 2: To Develop a Lightweight Trajectory Forecasting Module}

The second objective is to instill the system with a sense of "physics" and "object permanence."

\begin{itemize}
    \item \textbf{Rationale: } To survive occlusions and reduce sampling rates, the system must be able to "hallucinate"
          the future position of an object.
    \item \textbf{Implementation: } I will design and train a Recurrent Neural Network (RNN), specifically utilizing
          Long Short-Term Memory (\cite[LSTM]{hochreiter_long_1997}) or N-BEATS units. This model will take a
          lightweight input vector (Centroid $x, y$ Bounding Box $w, h$) rather than heavy image features. It will be trained on
          the YouTube-VOS trajectory dataset to learn typical motion patterns (linear, arc, acceleration). This module will act
          as the "autopilot" for the system between keyframes.
\end{itemize}

\subsection{Objective 3: To Adapt Geostatistical Kriging for Efficient Mask Reconstruction}

The third objective addresses the mathematical challenge of generating the mask shape without running a neural network.

\begin{itemize}
    \item \textbf{Rationale: } While the \cite[LSTM]{hochreiter_long_1997} predicts where the object is, it does
          not describe shape. Linear interpolation is too rigid for organic objects (e.g., a running dog).\
    \item \textbf{Implementation: } I will implement Universal Kriging and Radial Basis Function (RBF) interpolation.
          Crucially, to avoid the $O(N^3)$ computational complexity of Kriging, this objective includes the development of
          a "Contour-Vectorization" method. We will convert the binary raster masks from the keyframes into sparse polygon
          vertices, apply Kriging only to these vertices to model their deformation over time, and then rasterize the result.
          This transforms a pixel-heavy problem into a lightweight vector problem.
\end{itemize}

\subsection{Objective 4: To Validate the System via Comprehensive Benchmarking}

The final objective is to provide empirical evidence of the system's performance relative to the current State-of-the-Art
(\cite[SAM 2]{ravi_sam_2024} 2 / \cite[VideoLISA]{bai_one_2024}).

\begin{itemize}
    \item \textbf{Rationale: } Theoretical novelty is insufficient; the system must work on standard datasets to be academically valid.
    \item \textbf{Implementation: } The system will be tested on the MeViS (Motion Expressions Video Segmentation) dataset for reasoning
          capability and the DAVIS dataset for occlusion handling. I will measure:
          \begin{enumerate}
              \item \textbf{Semantic Accuracy: } Using the Jaccard Index $(J)$ and F-score $(F)$
              \item \textbf{Temporal Stability: } Using the Temporal Consistency Score $(T_stab)$
              \item \textbf{Efficiency: } Measuring FPS and Peak VRAM usage on a consumer-grade GPU (simulating edge constraints).
          \end{enumerate}
\end{itemize}